\documentclass[12pt]{article} 
\usepackage{natbib}
\usepackage{a4wide}
\usepackage{lineno}
\usepackage{cancel}
\usepackage{longtable}

\newif\ifbuch
\global\buchtrue
\global\buchfalse

\newif\ifdetail
\global\detailfalse
%\global\detailtrue

\usepackage{color}
\usepackage[pdftex]{graphicx}
\usepackage[pdftex,colorlinks,
                citecolor=darkgreen,linkcolor=darkblue,urlcolor=darkblue]{hyperref}
\definecolor{darkgreen}{rgb}{0.1,0.4,0.1}
\definecolor{darkblue}{rgb}{0.1,0.1,0.3}

\usepackage{amssymb}

\newcommand{\beq}  { \begin{eqnarray*} }
\newcommand{\eeq}  { \end{eqnarray*}}
\newcommand{\beeq}  { \begin{eqnarray*} }
\newcommand{\eeeq}  { \end{eqnarray*}}
\newcommand{\eq }  [1]{Eq.~(\ref{#1})}  % Eq. (#1)
\newcommand{\sect}  [1]{Sec.~(\ref{#1})}  % Eq. (#1)
\newcommand{\fig}  [1]{Fig.~\ref{#1}}   % Fig. #1
\renewcommand{\v}[1]{{\mbox{\boldmath$ #1 $}}}  % Vektor durch Fettdruck
\newcommand{\m}    [1]{{\bf  #1 }}              % Matrix in sans serif
\newcommand{\gz}[1] {{{d #1} \over {dz}}}
\newcommand{\pt}[1]{{\partial_t #1}}
\newcommand{\ptt}[1]{{{\partial^2 #1} \over {\partial t^2}}}
\newcommand{\pz}[1]{{\partial_z #1}}
\newcommand{\pzz}[1]{{\partial_{zz} #1}}
\newcommand{\px}[1]{{\partial_x #1}}
\newcommand{\py}[1]{{\partial_y #1}}
\newcommand{\p}{{\partial}}
\newcommand{\vn}{{\v \nabla}}
\newcommand{\tl}{\tilde}
\newcommand{\lk}{\left(}
\newcommand{\rk}{\right)}
\newcommand{\rvec} [1]{
    \raisebox{-1.5ex}{$\stackrel{\textstyle #1}{\neg}$} }% rotierter Vektor
\newcommand{\Vec}[1]{\left(\begin{array}{c} #1 \end{array}\right) }
\newcommand{\Bar}[1]{ \overline{ #1} }

\begin{document} 

\section*{Documentation for the simple non-hydrostatic model nh\_simple}

nh\_simple is a finite-difference sibling of the pseudo-spectral model ps3D:
it integrates the same scaled non-hydrostatic Boussinesq equations on the
same kind of domain (horizontally always periodic, vertically periodic or
bounded by rigid lids), but discretises them on an Arakawa C grid in
physical space and uses the Fourier transform only for the pressure
inversion. The LES closures are ported from ps3D. Within pyOM3, nh\_simple
is one of the ocean components that can be coupled to the spectral surface
wave model wim\_simple; the hooks of the Lagrangian-mean formulation of
that coupling are noted where they appear.

\subsection*{Continuous equations}

The scaled non-hydrostatic Boussinesq equations (see the scaling section
of the ps3D documentation for the derivation) are
\beq
 \pt{\v u } + f \rvec{\v u}  +   \vn p =   
  - Ro \lk \v u \cdot \vn \v u +  w \pz{} \v u \rk
 ~~,~~
    \pt{b}  + N^2  w    =  -Ro \lk  \v u \cdot \vn b  +  w \pz{} b \rk
 \\
   \delta^2 \pt w +   \pz p -  b =  - Ro \, \delta^2  \lk   \v u \cdot \vn w +  w \pz{} w \rk
\eeq
with the diagnostic relation $\vn \cdot \v u + \pz w = 0$. Vectors are
horizontal, $\rvec{\v u}$ denotes anticlockwise rotation by $90^o$,
$Ro$ is the Rossby number and $\delta = H/L$ the aspect ratio;
$f$ and the background stability frequency $N$ are constants ($f$-plane,
uniform stratification). The buoyancy $b$ is the deviation from the
background stratification $N^2 z/Ro$; the term $N^2 w$ is the advection of
this background. Kinetic and potential energy are
$\v u^2/2 + \delta^2 w^2/2$ and $b^2/2N^2$. Setting $Ro = 1$ and
$\delta^2 = 1$ recovers the plain dimensional equations, which is the
configuration required by the LES closures below.

\subsection*{Grid and spatial discretisation}

All fields live on an Arakawa C grid with uniform spacings
$\Delta x, \Delta y, \Delta z$: $b$ and $p$ at cell centres, $u,v,w$ on the
respective cell faces. The momentum advection is the second-order
energy-conserving flux form, e.g.\ for $u$
\beq
 \lk \pt u \rk_{adv} = -Ro \lk \delta_x^- \Bar{\Bar u^{x+} \Bar u^{x+}}
 + \delta_y^- \Bar{ \Bar u^{y+} \Bar v^{x+}}
 + \delta_z^- \Bar{ \Bar u^{z+} \Bar w^{x+}} \rk
\eeq
in the finite-difference and averaging operator notation of the ps3D
documentation ($\delta^\pm$ forward/backward differences, $\Bar{\cdot}^\pm$
two-point averages), and analogously for $v$ and $w$ (the $w$ tendency
carries the factor $\delta^2$). The Coriolis force uses the standard
four-point average of the C grid.

For the buoyancy advection three flux-form schemes are available, selected
by attribute: the second-order centred scheme (default, energy conserving),
the second-order scheme with Superbee flux limiter
(\verb|buoyancy_superbee_advection|, from MITgcm; monotone, and its
implicit scalar mixing can replace the explicit SGS buoyancy diffusion of
the LES closures, see below), and the third-order upwind scheme
(\verb|buoyancy_upwind3_advection|, from MITgcm). Momentum advection is
always the energy-conserving second-order form.

\subsection*{Pressure solver}

The velocity is first advanced with all terms except the pressure gradient
to an intermediate value $\v u^*, w^*$; the pressure then enforces the
divergence constraint,
\beq
 \lk \delta_x^- \delta_x^+ + \delta_y^- \delta_y^+
   + \delta_z^- \delta_z^+/\delta^2 \rk p =
 \frac{1}{\Delta t} \lk \delta_x^- u^* + \delta_y^- v^* + \delta_z^- w^* \rk
\eeq
followed by the projection
$\v u = \v u^* - \Delta t\, \vn^+ p$ and
$w = w^* - (\Delta t/\delta^2)\, \delta_z^+ p$.
As shown in the ps3D documentation, the Fourier transform of the discrete
one-dimensional operator $\delta^-\delta^+$ is
$-\hat k^+ \hat k^- = -2\lk 1 - \cos k \Delta \rk/\Delta^2$, which
approaches $-k^2$ for $k \Delta \to 0$; the discrete Poisson equation is
therefore diagonal in spectral space with the eigenvalue
$\hat k^2 = \hat k_x^+ \hat k_x^- + \hat k_y^+ \hat k_y^-
+ \hat k_z^+ \hat k_z^-/\delta^2$,
so the discrete divergence of the projected velocity vanishes to round-off.

Two cases are implemented, selected by \verb|periodic_in_z|. For the
triple-periodic domain a 3D FFT (parallel, through the p3dfft library)
inverts the full eigenvalue; the mean mode $\hat k^2 = 0$ is set to zero.
For rigid lids at top and bottom, a 2D FFT in $x,y$ is combined with a
tridiagonal solve in $z$ (Thomas algorithm) with Neumann boundary
conditions at the lids; the horizontal-mean column
($\hat k_x = \hat k_y = 0$), for which the tridiagonal system is singular,
is integrated directly by the second-order recurrence
$p_{k+1} = 2 p_k - p_{k-1} + \delta^2 \Delta z^2 f_k$ with $p_0 = 0$.

\subsection*{Boundary conditions}

The domain is always periodic in $x$ and $y$ (halo wrap or MPI exchange).
In $z$ it is either periodic (\verb|periodic_in_z| true) or bounded by
rigid free-slip lids: all vertical fluxes --- advective, frictional,
diffusive and the LES fluxes --- are set to zero on the top and bottom
faces, and $w$ vanishes there.

\subsection*{Friction and mixing}

Harmonic friction with constant coefficients acts on the momentum when the
attributes are present: with both \verb|Ah| and \verb|Av| the full operator
$A_h \vn_h^2 + A_v \p_{zz}$, with only one of them the corresponding
horizontal-only or vertical-only part. Harmonic buoyancy mixing requires
both \verb|Kh| and \verb|Kv|. Biharmonic friction and mixing
($-A_{hbi}\vn_h^4$ etc.) are activated by the pairs \verb|Ahbi|,
\verb|Avbi| for momentum and \verb|Khbi|, \verb|Kvbi| for buoyancy. All
these operators are applied with a forward Euler step outside the
Adams--Bashforth interpolation of the adiabatic tendencies, so their
coefficients are subject to the usual explicit diffusive stability limits.

\subsection*{Time stepping}

The adiabatic tendencies (advection, Coriolis, buoyancy force) are
interpolated with the multi-step Adams--Bashforth scheme,
\beq
(c_A, c_B, c_C) = (1,0,0) ~\to~ \lk \frac{3}{2}, -\frac{1}{2}, 0\rk
 ~\to~ \lk \frac{23}{12}, -\frac{16}{12}, \frac{5}{12} \rk
\eeq
i.e.\ a forward Euler start, one AB2 step and third-order Adams--Bashforth
from the third step on, with the tendencies held in a rotating three-level
buffer. Friction, mixing and the LES closures update the fields directly
with forward Euler before the Adams--Bashforth increment; the pressure
projection concludes the step. A user-defined forcing can be added to the
tendencies at the current time level through the optional method
\verb|set_forcing(tau)| (it then enters the Adams--Bashforth
interpolation), or applied as a direct increment through
\verb|set_forcing_direct()| after the Adams--Bashforth update and before
the pressure solve. The latter bypasses the AB weights and is the hook for
the Stokes tendency $\Delta \v u_S$ of the Lagrangian-mean ($\v u_L$)
formulation of the wave coupling, which must be a plain forward increment
and must precede the projection so that it acts on its divergence.

\subsection*{LES closures}

The LES closures of ps3D (\verb|les_closures.f90|) are ported as the mixin
class \verb|les| (module \verb|les.py|): the Smagorinsky closure, the
filtered-structure-function closure of Ducros et al.\ (1996), and a
Deardorff-type prognostic TKE closure, plus an optional Peclet bound on
the horizontal viscosities. The closures are formulated dimensionally and
refuse to run unless $Ro = 1$ and $\delta^2 = 1$. One call of
\verb|les_closure()| per time step (after the harmonic friction, before
the Adams--Bashforth increment, e.g.\ hooked through \verb|set_forcing|)
computes the strain-rate tensor on its natural C-grid positions,
evaluates the enabled closure(s), and applies the anisotropic viscous and
diffusive flux divergences to $u,v,w$ and $b$ with a forward Euler step.

\emph{Smagorinsky} (\verb|enable_smagorinsky|):
$\nu = (c_s \Delta)^2 \sqrt{2 S_{ij} S_{ij}}$ with
$\Delta = (\Delta x \Delta y \Delta z)^{1/3}$ and
$c_s = $ \verb|smag_coeff|; the diffusivity is $\nu/Pr$ with the turbulent
Prandtl number \verb|prandtl|. The viscosity is evaluated separately at
the four staggered positions where its fluxes are needed (cell centre,
horizontal corner, above $u$, above $v$).

\emph{Ducros et al.\ (1996)} (\verb|enable_ducros_96|): the viscosity is
built from the second-order structure function of the high-pass filtered
velocity (two iterated horizontal Laplace filters, three with
\verb|enable_three_level_filter|).

\emph{Prognostic TKE} (\verb|enable_tke_closure|): a Deardorff-type
closure with prognostic subgrid TKE on the $w$ grid, a
stratification-limited length scale
$\ell = \min(\Delta,\, 0.76 \sqrt{e/N^2})$, dissipation
$(0.19 + 0.51 \ell/\Delta)\, e^{3/2}/\ell$, TKE diffusion with
coefficient \verb|c_eke| times the eddy diffusivity, and shear and
buoyancy production diagnosed from the actual discrete mixing fluxes of
the previous step. Viscosity and diffusivity are
$\nu = c_m \ell \sqrt{e}$ and $\kappa = (1 + 2\ell/\Delta)\, \nu$ with
$c_m = $ \verb|c_m|; $N^2$ always includes the background $N_0^2$.

\emph{Stability caps.} All closures cap their viscosities and
diffusivities against the explicit diffusive limit of the forward Euler
update. The caps are direction-aware: each staggered viscosity is capped
with the most restrictive grid spacing among the directions its fluxes
act in (e.g.\ the cell-centre viscosity by $\min(\Delta x, \Delta y,
\Delta z)^2/(4\Delta t)$). This deviates deliberately from the original
ps3D caps, which use geometric means of the spacings and are valid only
for near-isotropic grids: for $\Delta x \gg \Delta z$ they admit
vertically unstable viscosities.

\emph{Boundary treatment.} With rigid lids the vertical shear seen by the
closures is zeroed on the lids (\verb|les_freeslip_strain|, default on),
consistent with the free-slip zero-flux condition of the mixing;
otherwise the zero halo above the lid produces a spurious shear that
inflates the viscosity in the top and bottom cells.

\emph{Scalar mixing.} The explicit SGS buoyancy diffusion can be switched
off (\verb|les_mix_b|) when a flux-limited buoyancy advection scheme
(Superbee) provides the scalar SGS mixing implicitly; under the TKE
closure the resolved buoyancy production is then zero.

\emph{Peclet bound} (\verb|enable_peclet_bound|): where the grid Peclet
number of the resolved flow exceeds \verb|Peclet_limit|, the horizontal
viscosities are raised proportionally. The bound uses the prognostic
(advecting) velocity, i.e.\ $\v u_L$ in the Lagrangian-mean formulation.

\emph{Coupled wave--ocean setups.} When the prognostic velocity is the
Lagrangian mean $\v u_L$, the attribute \verb|les_mix_eulerian| makes the
strains and mixing fluxes act on the Eulerian velocity
$\v u = \v u_L - \v u_S$ (attributes \verb|u_S|, \verb|v_S| and optionally
\verb|w_S| supplied by the coupler), with the tendency added to the
prognostic $\v u_L$ --- the down-gradient closure parameterises Eulerian
fluctuations, so mixing $\v u_L$ would spuriously diffuse the Stokes
shear. This is the LES counterpart of the \verb|mix_eulerian| option of
the TKE closure in pe\_simple.

\subsection*{Domain decomposition}

The parallelisation follows the pencil decomposition of the p3dfft
library used by the pressure solver: the $x$ direction is kept local on
every processor, and the constructor arguments \verb|n_pes_x| and
\verb|n_pes_y| define the processor grid over the $y$ and $z$ directions
(internally the blocks are addressed as $y$/$z$ blocks). Model arrays are
Fortran-ordered $(N_x, N_y, N_z)$ with halo points on all sides; the
pressure solver requires this ordering. Halo exchanges wrap periodically
in $x$ and $y$; in $z$ they wrap only for the triple-periodic case.

\subsection*{Snapshot diagnostics}

Two snapshot classes live in the module \verb|snapshot.py|; either
replaces the model class as the base of the configuration and overrides
\verb|setup| and \verb|diagnose|. Both read the attribute \verb|snapint|
(snapshot interval in seconds, required), take the setup arguments
\verb|snap_file_name| and \verb|show_each_time_step| (default on: print
the model time every step; off: only at output steps), and write at
$n = 0$ and whenever $n \bmod (\mathrm{snapint}/\Delta t) = 0$. Both
record the variables $u,v,w,b,p$ on their staggered coordinate axes
(\verb|xt|/\verb|yt|/\verb|zt| for cell centres and
\verb|xu|/\verb|yu|/\verb|zu| for the faces, shifted by half a spacing);
the vertical axis is negative downwards with the surface at $z = 0$ and
the cell centres at $z = -\Delta z/2, \ldots$. Halo points are stripped
before writing.

\emph{snap\_nc4} writes a single netCDF file (default \verb|'snap.cdf'|).
At start-up it detects whether the netCDF4/HDF5 build supports parallel
I/O (\verb|__has_parallel4_support__| or pnetcdf) and reports the result.
In parallel mode the file is created and opened collectively by all PEs,
and every PE writes its subdomain into the global arrays in a single
collective pass per snapshot. Without parallel support the original
serialised path is used: PE0 creates the file and extends the time
dimension, then each PE in turn opens the file and writes its block,
synchronised with barriers --- correct on any build, but increasingly
slow with many PEs. The internal Fortran ordering is transposed on the
fly, so the file layout (time, $z$, $y$, $x$) is the same in both modes.

\emph{snap\_zarr\_gather} writes a single zarr store (default
\verb|'snap.zarr'|) with field arrays of shape (time, $n_x$, $n_y$,
$n_z$), always $x$-first regardless of the internal ordering, one time
slice per chunk. The design goal is a Lustre-friendly file count: with
the default \verb|z_chunks| $= 1$ each snapshot adds one chunk file per
field --- of the order of six files per snapshot regardless of the number
of PEs --- instead of the one-file-per-PE-per-field fragmentation of a
per-block chunking, which can overwhelm the metadata server of a parallel
file system at large PE counts. Every PE participates in the collective
gather of each field (\verb|gather_array_mpi|), but only PE0 touches the
store, writing a full horizontal slab of $\lceil n_z/\verb|z_chunks|
\rceil$ levels per chunk; increase \verb|z_chunks| only if a single field
slab is too large to hold and write in one piece. Fields are gathered and
written one at a time to bound the PE0 memory to one global field. The
time axis and the field arrays are resized (appended) per snapshot, so
the store remains valid after an interrupted run; an existing store of
the same name is removed at start-up. The on-disk format is pinned to
zarr v2 even when a zarr-python v3 library writes it, so the output stays
readable by v2 environments; the cosmetic consolidated-metadata warning
of v3 is silenced. Compared to the netCDF path this trades the gather
(one global field on PE0) for independence from parallel-I/O library
support.

\section*{Model parameters and options}

All parameters and options are attributes of the model instance, set
either through the parameter dictionary passed at construction or by
assigning the attribute before the time loop; options with a default may
be left unset.

\subsection*{Required parameters (no default)}

\begin{longtable}{@{}llp{8.4cm}@{}}
\hline
Attribute & Unit & Meaning \\
\hline
\endhead
\verb|nx|, \verb|ny|, \verb|nz| & --- & global grid size (cells, without
  halo points) \\
\verb|dx|, \verb|dy|, \verb|dz| & m & grid spacing (uniform) \\
\verb|dt| & s & time step \\
\verb|f0| & 1/s & Coriolis parameter ($f$-plane) \\
\verb|N0| & 1/s & background stability frequency \\
\verb|Ro| & --- & Rossby number of the scaled equations ($1$ for the
  dimensional form; the LES closures require $1$) \\
\verb|dsqr| & --- & squared aspect ratio $\delta^2$ ($1$ for the fully
  non-hydrostatic dimensional form; the LES closures require $1$) \\
\verb|periodic_in_z| & --- & triple-periodic domain if true; rigid
  free-slip lids at top and bottom if false. $x$ and $y$ are always
  periodic \\
\hline
\end{longtable}

\noindent
The constructor additionally takes \verb|n_pes_x|, \verb|n_pes_y|
(processor grid over the $y$ and $z$ pencils, default $1 \times 1$;
their product must equal the communicator size), \verb|halo_points|
(default 1) and \verb|comm_in| (default: MPI world communicator).

\subsection*{Forcing hooks and coupling fields}

\begin{longtable}{@{}lp{10.6cm}@{}}
\hline
Attribute / method & Meaning \\
\hline
\endhead
\verb|set_forcing(tau)| & optional user method; adds forcing to the
  tendencies at time level \verb|tau| (enters the Adams--Bashforth
  interpolation). Also the natural hook to call \verb|les_closure()| \\
\verb|set_forcing_direct()| & optional user method; direct increment after
  the Adams--Bashforth update and before the pressure solve, bypassing the
  AB weights. Used for the Stokes tendency $\Delta \v u_S$ of the
  Lagrangian-mean ($\v u_L$) formulation \\
\verb|u_S|, \verb|v_S|, \verb|w_S| & Stokes velocity from the wave model
  ($w_S$ optional); required by \verb|les_mix_eulerian| \\
\hline
\end{longtable}

\subsection*{Options with defaults}

\begin{longtable}{@{}llp{7.9cm}@{}}
\hline
Attribute & Default & Meaning \\
\hline
\endhead
\verb|buoyancy_superbee_advection| & off & buoyancy advection with the
  Superbee flux limiter instead of the centred second-order scheme \\
\verb|buoyancy_upwind3_advection| & off & third-order upwind buoyancy
  advection \\
\verb|Ah|, \verb|Av| & unset & harmonic viscosities (m$^2$/s); either one
  alone activates the horizontal-only / vertical-only operator \\
\verb|Kh|, \verb|Kv| & unset & harmonic buoyancy diffusivities (m$^2$/s);
  both must be present \\
\verb|Ahbi|, \verb|Avbi| & unset & biharmonic viscosities (m$^4$/s); both
  must be present \\
\verb|Khbi|, \verb|Kvbi| & unset & biharmonic buoyancy diffusivities
  (m$^4$/s); both must be present \\
\hline
\end{longtable}

\subsection*{LES options (read at \texttt{init\_les})}

\begin{longtable}{@{}llp{7.9cm}@{}}
\hline
Attribute & Default & Meaning \\
\hline
\endhead
\verb|enable_smagorinsky| & off & Smagorinsky closure \\
\verb|enable_ducros_96| & off & Ducros et al.\ (1996)
  structure-function closure \\
\verb|enable_tke_closure| & off & Deardorff-type prognostic TKE closure \\
\verb|enable_three_level_filter| & off & three instead of two iterated
  Laplace filters (Ducros closure) \\
\verb|enable_peclet_bound| & off & raise horizontal viscosities where the
  grid Peclet number exceeds \verb|Peclet_limit| \\
\verb|les_mix_eulerian| & off & strains and mixing act on the Eulerian
  velocity $\v u_L - \v u_S$ ($\v u_L$ formulation) \\
\verb|les_freeslip_strain| & on & zero the vertical shear seen by the
  closure at the rigid lids \\
\verb|les_mix_b| & on & explicit SGS buoyancy diffusion; switch off when a
  flux-limited advection scheme provides the scalar SGS mixing \\
\verb|smag_coeff| & $0.148$ & Smagorinsky coefficient \\
\verb|prandtl| & $0.47$ & turbulent Prandtl number \\
\verb|min_smag| & $0$ & viscosity floor (m$^2$/s) \\
\verb|Peclet_limit| & $50$ & grid Peclet number bound \\
\verb|c_m| & $0.1$ & TKE-closure viscosity constant \\
\verb|c_eke| & $2$ & TKE diffusion constant \\
\verb|tke_min| & $10^{-8}$ & TKE floor (m$^2$/s$^2$) \\
\verb|les_epsln| & $10^{-20}$ & numerical floor of the closure \\
\hline
\end{longtable}

\subsection*{Snapshot diagnostics}

Both snapshot classes (module \verb|snapshot.py|) read the attribute
\verb|snapint| (snapshot interval in seconds, required) and take the setup
arguments \verb|snap_file_name| (defaults \verb|'snap.cdf'| /
\verb|'snap.zarr'|) and \verb|show_each_time_step| (default on);
\verb|snap_zarr_gather| additionally takes \verb|z_chunks| (default 1:
one chunk file per field and snapshot; larger values split the vertical
into $\lceil n_z/\verb|z_chunks| \rceil$-level chunks).

\end{document}
